% 1-WL 局限：六环 C6 与两个不相交三角形；初始同标记 A，每点每轮看到两个同类邻居，签名恒为 (A,{{A,A}})。
\begin{tikzpicture}[font=\figurefont\fontsize{8}{10}\selectfont,
  vtx/.style={circle,draw=FigureStroke,line width=.8pt,minimum size=6.2mm,inner sep=0pt,text=FigureInk,fill=FigureBlueFill},
  edge/.style={draw=FigureStroke,line width=.85pt},
  note/.style={font=\FigureNoteFont\fontsize{7.5}{9}\selectfont,text=FigureMuted,align=center}]
\begin{scope}
  \foreach \i in {1,...,6} \node[vtx] (a\i) at ({90-60*(\i-1)}:1.05) {$A$};
  \foreach \i/\j in {1/2,2/3,3/4,4/5,5/6,6/1} \draw[edge] (a\i)--(a\j);
  \node at (0,-1.8) {(a) 六环：1 个分量};
\end{scope}
\begin{scope}[xshift=3.3cm]
  \foreach \x/\p in {0/b,2.25/c}{
    \node[vtx] (\p1) at (\x,0.75) {$A$};
    \node[vtx] (\p2) at (\x-0.65,-0.4) {$A$};
    \node[vtx] (\p3) at (\x+0.65,-0.4) {$A$};
    \draw[edge] (\p1)--(\p2)--(\p3)--(\p1);
  }
  \node at (1.12,-1.8) {(b) 两个三角形：2 个分量};
\end{scope}
\node[note] at (2.9,-2.45) {两图每个顶点的签名都是 $(A,\{\!\{A,A\}\!\})$，各轮类别计数相同};
\end{tikzpicture}
